%!TEX root = ../../main.tex
\section{M-RQ4에 대한 답}
\label{sec:val-answer-rq4}

예측의 이득은 초기 우세와 시간의 유연한 함수를 넘어서 존재한다. 균형 구간 안에서는 더 작고(T에서는 구간의 상단이 0에 가깝다), 작은 변화를 제외해도 유지되며, 점수 분해상 판별 성분의 증가와 함께 나타난다. 관측 갱신으로 층화하면 새 프레임이 들어온 셀에서 전체와 같은 부호이고 같은 프레임 셀에서는 0 부근이다(C26). 이후 패치와 다른 지역에서는, 평가한 외부 표본에서 한타의 우세가 관측되지 않았다. 지역별로 보면 KR에서는 차이가 불확실하고 NA1에서는 기준선에 불리한 차이가 95\%와 98.75\% 구간 모두에서 관측되었다. 소규모 교전의 우세는 두 표본에서 95\% 구간 수준에서 관측되었고 가족 보정 수준에서는 불확실하다(C27). 이 결과는 동결된 평가기, 교전 정의, 평가된 절차와 표본에 조건부이며, 공통 해석 범위는 제 \ref{sec:disc-scope} 절에 둔다.
